
\documentclass{article}
\usepackage{colortbl}
\usepackage{makecell}
\usepackage{multirow}
\usepackage{supertabular}

\begin{document}

\newcounter{utterance}

\centering \large Interaction Transcript for game `cladder', experiment `full\_v1.5\_default', episode 4010 with Qwen3.5{-}9B.
\vspace{24pt}

{ \footnotesize  \setcounter{utterance}{1}
\setlength{\tabcolsep}{0pt}
\begin{supertabular}{c@{$\;$}|p{.15\linewidth}@{}p{.15\linewidth}p{.15\linewidth}p{.15\linewidth}p{.15\linewidth}p{.15\linewidth}}
    \# & \multicolumn{2}{c}{Player} && \multicolumn{2}{c}{Game Master} \\
    \hline

    \theutterance \stepcounter{utterance}  
    & & & \multicolumn{4}{p{0.6\linewidth}}{
        \cellcolor[rgb]{0.9,0.9,0.9}{
            \makecell[{{p{\linewidth}}}]{
                \texttt{\tiny{[P1$\langle$GM]}}
                \texttt{You are an expert in causal inference. The following question is not a typical commonsense query, but rather a meticulously designed question created by a professor specializing in causal inference, intended to assess the students' mastery of the course content.} \\
\texttt{Imagine a self{-}contained, hypothetical world with only the following conditions, and without any unmentioned factors or causal relationships: Liking spicy food has a direct effect on parents' social status and child's intelligence. Other unobserved factors has a direct effect on parents' social status and child's intelligence. Parents' social status has a direct effect on child's intelligence. Other unobserved factors is unobserved. The overall probability of liking spicy food is 51\%. The probability of not liking spicy food and intelligent child is 18\%. The probability of liking spicy food and intelligent child is 34\%. Is the chance of intelligent child larger when observing liking spicy food?} \\
\texttt{Start your answer with 'Yes' or 'No', followed by additional reasoning or evidence to support your explanation.} \\
            }
        }
    }
     \\ \\

    \theutterance \stepcounter{utterance}  
    & \multicolumn{4}{p{0.6\linewidth}}{
        \cellcolor[rgb]{0.561,0.737,0.561}{
            \makecell[{{p{\linewidth}}}]{
                \texttt{\tiny{[P1$\rangle$GM]}}
                \texttt{Yes} \\
\\ 
\texttt{**Reasoning:**} \\
\\ 
\texttt{To determine if the chance of an intelligent child is larger when observing that parents like spicy food, we must compare the conditional probability \$P(\textbackslash{}text\{Intelligent\} \textbackslash{}mid \textbackslash{}text\{Spicy\})\$ with the marginal probability \$P(\textbackslash{}text\{Intelligent\})\$.} \\
\\ 
\texttt{1.  **Calculate the marginal probability of an intelligent child, \$P(\textbackslash{}text\{Intelligent\})\$:**} \\
\texttt{    The problem provides the joint probabilities for the two mutually exclusive states of "Liking Spicy Food" (Spicy) and "Not Liking Spicy Food" (Not Spicy).} \\
\texttt{    *   \$P(\textbackslash{}text\{Spicy\} \textbackslash{}cap \textbackslash{}text\{Intelligent\}) = 34\textbackslash{}\%\$} \\
\texttt{    *   \$P(\textbackslash{}text\{Not Spicy\} \textbackslash{}cap \textbackslash{}text\{Intelligent\}) = 18\textbackslash{}\%\$} \\
\\ 
\texttt{    Summing these gives the total probability of having an intelligent child:} \\
\texttt{    \$\$P(\textbackslash{}text\{Intelligent\}) = 34\textbackslash{}\% + 18\textbackslash{}\% = 52\textbackslash{}\%\$\$} \\
\\ 
\texttt{2.  **Calculate the conditional probability of an intelligent child given liking spicy food, \$P(\textbackslash{}text\{Intelligent\} \textbackslash{}mid \textbackslash{}text\{Spicy\})\$:**} \\
\texttt{    Using the definition of conditional probability \$P(A \textbackslash{}mid B) = \textbackslash{}frac\{P(A \textbackslash{}cap B)\}\{P(B)\}\$:} \\
\texttt{    *   We know \$P(\textbackslash{}text\{Spicy\} \textbackslash{}cap \textbackslash{}text\{Intelligent\}) = 34\textbackslash{}\%\$.} \\
\texttt{    *   We are given the overall probability of liking spicy food, \$P(\textbackslash{}text\{Spicy\}) = 51\textbackslash{}\%\$.} \\
\\ 
\texttt{    \$\$P(\textbackslash{}text\{Intelligent\} \textbackslash{}mid \textbackslash{}text\{Spicy\}) = \textbackslash{}frac\{34\textbackslash{}\%\}\{51\textbackslash{}\%\} \textbackslash{}approx 0.6667 \textbackslash{}text\{ or \} 66.7\textbackslash{}\%\$\$} \\
\\ 
\texttt{3.  **Comparison:**} \\
\texttt{    *   \$P(\textbackslash{}text\{Intelligent\} \textbackslash{}mid \textbackslash{}text\{Spicy\}) \textbackslash{}approx 66.7\textbackslash{}\%\$} \\
\texttt{    *   \$P(\textbackslash{}text\{Intelligent\}) = 52\textbackslash{}\%\$} \\
\\ 
\texttt{    Since \$66.7\textbackslash{}\% > 52\textbackslash{}\%\$, the probability of having an intelligent child is indeed larger when observing that the parents like spicy food.} \\
\\ 
\texttt{*Note on Causal Structure:* While the problem describes a causal graph where "Liking Spicy Food" affects "Child's Intelligence" (directly) and "Parents' Social Status" (which also affects "Child's Intelligence"), the question asks for a probabilistic comparison based on the observed data provided. The calculation above relies strictly on the numerical constraints given (\$P(\textbackslash{}text\{Spicy\})\$, \$P(\textbackslash{}text\{Spicy\} \textbackslash{}cap \textbackslash{}text\{Intelligent\})\$, and} \\
            }
        }
    }
    & & \\ \\

    \theutterance \stepcounter{utterance}  
    & & & \multicolumn{2}{p{0.3\linewidth}}{
        \cellcolor[rgb]{0.9,0.9,0.9}{
            \makecell[{{p{\linewidth}}}]{
                \texttt{\tiny{[GM$|$GM]}}
                \texttt{Yes} \\
            }
        }
    }
    & & \\ \\

    \theutterance \stepcounter{utterance}  
    & & & \multicolumn{2}{p{0.3\linewidth}}{
        \cellcolor[rgb]{0.9,0.9,0.9}{
            \makecell[{{p{\linewidth}}}]{
                \texttt{\tiny{[GM$|$GM]}}
                \texttt{yes} \\
            }
        }
    }
    & & \\ \\

    \theutterance \stepcounter{utterance}  
    & & & \multicolumn{2}{p{0.3\linewidth}}{
        \cellcolor[rgb]{0.9,0.9,0.9}{
            \makecell[{{p{\linewidth}}}]{
                \texttt{\tiny{[GM$|$GM]}}
                \texttt{game\_result = WIN} \\
            }
        }
    }
    & & \\ \\

\end{supertabular}
}

\end{document}
